\section{Synthesis: The Teacher and the Trial}
\label{sec:origen-synthesis}

Origen's life has an intelligible center. A child formed by Scripture and his
father's execution made reading an act of total discipleship. He built schools
in which grammar and philosophy served the biblical Word, a manuscript workshop
in which textual difference could be inspected, and a preaching practice in
which the same text judged ordinary Christian desire. When imperial power
required sacrifice, the theory of confession became bodily fact.

The center did not prevent overreach. Origen wanted divine justice, freedom,
the diversity of creatures, embodiment, and the final victory of goodness to
form one coherent economy. Where apostolic proclamation did not specify the
manner, he explored. Some explorations gave later theology indispensable
language: the Son's eternal generation, Scripture's spiritual depth, freedom
against fixed natures, and punishment ordered toward correction. Others tied
human embodiment to pre-existent fault, rendered resurrection vulnerable to
dissolution, or made the end more schematic than the Church could receive.

The difference between Origen and Origenism is therefore neither absolute nor
imaginary. Later ascetics, translators, accusers, and imperial theologians
selected and transformed his proposals. They must not be quoted as his direct
voice. Yet a founder is historically related to the possibilities his system
opens even when followers construct a doctrine he never stated in that form.
Reception history is fairest when it traces that relation instead of choosing
between total blame and total exoneration.

His Christian character should be narrated with the same restraint. The
sources establish disciplined prayer, radical generosity, service to churches,
formation of pupils, and fidelity under torture. They also establish an
institutional rupture, severe bodily practice, damaging polemical structures,
and disputed doctrine. No competent act inspected for this edition canonizes
him, and this project does not create one. Sanctity is not proved or disproved
by historical aggregation; what history can do is refuse a fictive cult while
making the cost and fruit of discipleship visible.

Three absences remain decisive. The Alexandrian proceedings are lost, so the
living Church's first case against Origen cannot be retried from full evidence.
Most of the Greek corpus is lost or mediated, so a perfectly reconstructed
``true Origen'' remains an aspiration. The death and burial have no continuous
near-contemporary narrative, so the confessor's last act cannot be made to say
what his books say.

Those absences do not exhaust the legacy. A durable gift is a
practice of connected attention: compare the wording; listen to interlocutors;
read within the Church's proclamation; bring intellect into prayer; test an
interpretation by the life it forms. His history adds the necessary reverse
tests: identify who paid and copied; ask whom an interpretation displaces;
distinguish inquiry from confession; and do not let a brilliant connection
outrun evidence or received faith.

The reception of this teacher is accordingly not resolved by simple approval
or by repeating an anathema. It continues wherever Christians read him with exact
gratitude---deeply enough to receive what is true, historically enough to know
whose words survive, and ecclesially enough to name the point at which
exploration can no longer govern belief.
